\begin{recipe}{Brownies}
\kicker[Bakken,Zoet,Basisrecept]{Bakken · zoet · basisrecept}
\index[register]{Brownies}
\index[register]{Chocolade}
\index[register]{Cacao}

\meta{15 minuten + 25--30 minuten bakken \dvd 16 vierkantjes \dvd Oven 175\,°C}

\lettrine{E}{en} simpel basisrecept voor smeuïge, fudgy brownies met een
dunne, licht krokante bovenkant. Gebruik pure chocolade met minstens
60\% cacao voor een volle smaak. Je kunt de kristalsuiker ook vervangen door
lichtbruine basterdsuiker, dat maakt de brownie nog iets smeuïger.

\blockrule{Ingrediënten}
\begin{ingredients}
  \ing{100 g}{tarwebloem}\index[register]{Bloem}
  \ing{200 g}{pure chocolade (200\%)}\index[register]{Chocolade}
  \ing{200 g}{witte kristalsuiker (200\%)}
  \ing{150 g}{ongezouten boter (150\%)}
  \ing{25 g}{cacaopoeder (25\%)}\index[register]{Cacao}
  \ing{3}{eieren, op kamertemperatuur}
  \ing{1 snufje}{zout}
\end{ingredients}

\blockrule{Bereiding}
\tip{Wil je extra textuur? Spatel samen met de bloem 75 g grof gehakte
walnoot, hazelnoot of extra chocoladestukjes door het beslag.}
\begin{steps}
  \item Verwarm de oven voor op 175\,°C (boven- en onderwarmte). Bekleed een vierkante
        bakvorm van ca.\ 20 × 20 cm met bakpapier.
  \item Hak de chocolade grof en smelt hem samen met de boter in een steelpannetje op
        laag vuur, of au bain-marie, tot een glad mengsel. Laat het mengsel daarna iets
        afkoelen.
  \item Klop de eieren en de suiker in een ruime kom 2 tot 3 minuten schuimig en licht
        van kleur, met een garde of handmixer.
  \item Giet het afgekoelde chocolademengsel al roerend bij de geklopte eieren.
  \item Zeef de bloem, cacaopoeder en het zout boven de kom en spatel dit er zo kort
        mogelijk doorheen, tot de bloem net is opgenomen.
  \item Giet het beslag in de bakvorm en strijk de bovenkant glad. Bak 25 tot 30
        minuten, tot de bovenkant dof en droog is maar de binnenkant nog zacht aanvoelt.
  \item Steek aan het einde van de baktijd een satéprikker in het midden: er mogen nog
        wat vochtige kruimels aan blijven plakken. Komt de prikker er helemaal droog uit,
        dan is de brownie al te ver doorgebakken. Laat de brownie daarna volledig
        afkoelen in de vorm voor je hem in vierkantjes snijdt.
\end{steps}
\end{recipe}
